\chapter{Problems}
\begin{quote}
  There is light at the end of the tunnel, but it is moving away at speed $c$.
\end{quote}


Here, we provide some problems for the reader to attempt.  Some are easier than others.  Infinity \cite{infinity} has a longer list than is included here; see \url{http://www.bmoc.maths.org/home/areals.pdf} for the relevant excerpt.
%\newcommand{\sollink}[1]{\footnote{A rough outline of the solution can be found on the AoPS thread \href{#1}{here}.}}
\newcommand{\sollink}[1]{}

For practice ``in real life'', i.e. determining when barycentric coordinates are feasible for a problem (here, they are), you can often simply take an arbitrary geometry problem that begins with the words ``triangle $ABC$'', and see if it's possible to get a barycentric solution.  Some will, of course, take much longer than others, but a good portion will have a reasonable solution.

\section{Problems}

\begin{enumerate}
  \item Prove Stewart's Theorem.

  \item Prove Routh's Theorem.

  \item (USA TST 2003/2) Let $ABC$ be a triangle and let $P$ be a point in its interior. Lines $PA$, $PB$, $PC$ intersect sides $BC$, $CA$, $AB$ at $D$, $E$, $F$, respectively. Prove that
\[ [PAF]+[PBD]+[PCE]=\frac{1}{2}[ABC]  \]
if and only if $P$ lies on at least one of the medians of triangle $ABC$. (Here $[XYZ]$ denotes the area of triangle $XYZ$.)

  \item (Mongolia TST 2000/6) In a triangle $ABC$, the angle bisector at $A$,$B$,$C$ meet the opposite sides at $A_1$,$B_1$,$C_1$, respectively. Prove that if the quadrilateral $BA_1B_1C_1$ is cyclic, then
  \[ \frac{AC}{AB+BC} = \frac{AB}{AC+BC} + \frac{BC}{BA+AC} \]

  \item (ISL 2005/G1) Given a triangle $ABC$ satisfying $AC+BC=3\cdot AB$. The incircle of triangle $ABC$ has center $I$ and touches the sides $BC$ and $CA$ at the points $D$ and $E$, respectively. Let $K$ and $L$ be the reflections of the points $D$ and $E$ with respect to $I$. Prove that the points $A$, $B$, $K$, $L$ lie on one circle.  \sollink{http://www.artofproblemsolving.com/Forum/viewtopic.php?p=2351823\#p2351823}

  \item (APMO 2005/5) In a triangle $ABC$, points $M$ and $N$ are on sides $AB$ and $AC$, respectively, such that $MB = BC = CN$. Let $R$ and $r$ denote the circumradius and the inradius of the triangle $ABC$, respectively. Express the ratio $MN/BC$ in terms of $R$ and $r$.  \sollink{http://www.artofproblemsolving.com/Forum/viewtopic.php?f=47\&t=30955\&p=2376234\#p2376234}

  \item (ISL 2008/G4) In an acute triangle $ ABC$ segments $ BE$ and $ CF$ are altitudes. Two circles passing through the point $ A$ and $ F$ and tangent to the line $ BC$ at the points $ P$ and $ Q$ so that $ B$ lies between $ C$ and $ Q$. Prove that lines $ PE$ and $ QF$ intersect on the circumcircle of triangle $ AEF$.

  \item (ISL 2001/G6) Let $ABC$ be a triangle and $P$ an exterior point in the plane of the triangle. Suppose the lines $AP$, $BP$, $CP$ meet the sides $BC$, $CA$, $AB$ (or extensions thereof) in $D$, $E$, $F$, respectively. Suppose further that the areas of triangles $PBD$, $PCE$, $PAF$ are all equal. Prove that each of these areas is equal to the area of triangle $ABC$ itself.  \sollink{http://www.artofproblemsolving.com/Forum/viewtopic.php?p=119203\&sid=f14625806c6d8ca5231f857e4519121e\#p119203}

  \item (IMO 2007/2) Consider five points $A$, $B$, $C$, $D$ and $E$ such that $ABCD$ is a parallelogram and $BCED$ is a cyclic quadrilateral. Let $\ell$ be a line passing through $A$. Suppose that $\ell$ intersects the interior of the segment $DC$ at $F$ and intersects line $BC$ at $G$. Suppose also that $EF = EG = EC$. Prove that $\ell$ is the bisector of angle $DAB$. \sollink{http://www.artofproblemsolving.com/Forum/viewtopic.php?f=451\&t=159837\&p=2660858\#p2660858}
\end{enumerate}


\section{Hints}
Hints have been obfuscated by ROT13.
\begin{enumerate}
  \item ABG UNEQ
  \item FGENVTUGSBEJNEQ
  \item HFR NERN SBEZ
  \item PVEPYR SBEZ FBYIR HIJ
  \item PVEPF GURA J VF ZVAHF N O
  \item QVFG SNPGBE GURA TRB VQ
  \item CBJRE CBVAG FUBJF NRCD PLPYVP GURA NATYR PUNFR
  \item PNERSHY FVTARQ NERNF
  \item ERS OPQ... SVAQ S T FLF RDA RYVZ P M SNPGBE
\end{enumerate}
